%! title: GRPO 与 PPO：Rollout、Scoring 与参数更新
%! date: 2026-10-26
%! slug: rl-grpo
%! series: 强化学习
%! number: 3
%! tags: 强化学习, PPO, GRPO, Rollout, 重要性采样
%! description: 把一次 PPO 迭代拆成 Rollout 与 Scoring 两次模型调用：一个串行、随机、慢，一个并行、确定、快；再说明两种 log 概率如何变成重要性权重 r_t。
\documentclass[UTF8,fontset=fandol]{ctexart}

% 版式
\usepackage{geometry}
\geometry{
  a4paper,
  margin=2.5cm,
  headheight=14pt,
}

% 主题色（克制、印刷友好）
\usepackage{xcolor}
\definecolor{themeaccent}{HTML}{1a5f7a}
\definecolor{themecolorthm}{HTML}{1a5f7a}
\definecolor{themecolordef}{HTML}{2d6a4f}
\definecolor{themecolorremark}{HTML}{8b6914}

% 定理块底纹浓度（xcolor：数字越大底色越深；旧版约 !4–!6，此处加深便于对比）
\newcommand{\thmshadecolor}{themecolorthm!18}
\newcommand{\defshadecolor}{themecolordef!14}
\newcommand{\remarkshadecolor}{themecolorremark!12}

% 数学
\usepackage{amsmath,amssymb}
\usepackage{mathtools}
\usepackage{amsthm}
\usepackage[anycase]{thmtools}
\usepackage{bm}

% 行距（可读性）
\usepackage{setspace}
\onehalfspacing

% 页眉页脚
\usepackage{fancyhdr}
\pagestyle{fancy}
\fancyhf{}
\fancyhead[L]{\small\color{themeaccent!70!black}\nouppercase{\leftmark}}
\fancyhead[R]{\small\thepage}
\renewcommand{\headrulewidth}{0.4pt}
\renewcommand{\headrule}{\color{themeaccent!50!black}\hrule height 0.4pt}

% 章节标题着色 + 装饰线（ctex）
\ctexset{
  section/format     = \Large\bfseries\color{themeaccent},
  section/aftername  = \hspace{0.5em},
  section/aftertitle = \par\vspace{-4pt}%
    {\color{themeaccent!40}\rule{\linewidth}{0.5pt}}\vspace{4pt},
  subsection/format  = \large\bfseries\color{themeaccent!88!black},
}

% 常用配对符号（mathtools）
\DeclarePairedDelimiter{\paren}{(}{)}
\DeclarePairedDelimiter{\sqbrk}{[}{]} % 勿用 \brack，与内核冲突
\DeclarePairedDelimiter{\curlbrace}{\{}{\}} % 勿用 \brace，与内核冲突
\DeclarePairedDelimiter{\abs}{\lvert}{\rvert}
\DeclarePairedDelimiter{\norm}{\lVert}{\rVert}

% 定理样式：彩色标题（底纹在 \declaretheorem 中用 shaded 指定）
\declaretheoremstyle[
  spaceabove=6pt,
  spacebelow=6pt,
  headfont=\bfseries\color{themecolorthm},
  notefont=\bfseries\color{themecolorthm!80!black},
  bodyfont=\itshape,
  headpunct={.},
  postheadspace=0.5em,
]{styletheorem}

\declaretheoremstyle[
  spaceabove=6pt,
  spacebelow=6pt,
  headfont=\bfseries\color{themecolordef},
  notefont=\bfseries\color{themecolordef!80!black},
  bodyfont=\normalfont,
  headpunct={.},
  postheadspace=0.5em,
]{styledefinition}

\declaretheoremstyle[
  spaceabove=6pt,
  spacebelow=6pt,
  headfont=\bfseries\color{themecolorremark},
  notefont=\bfseries\color{themecolorremark!80!black},
  bodyfont=\normalfont,
  headpunct={.},
  postheadspace=0.5em,
]{styleremark}

\declaretheorem[
  style=styletheorem,
  name=定理,
  numberwithin=section,
  shaded={bgcolor=\thmshadecolor, margin=0pt},
]{theorem}

\declaretheorem[
  style=styletheorem,
  name=引理,
  sibling=theorem,
  shaded={bgcolor=\thmshadecolor, margin=0pt},
]{lemma}

\declaretheorem[
  style=styletheorem,
  name=命题,
  sibling=theorem,
  shaded={bgcolor=\thmshadecolor, margin=0pt},
]{proposition}

\declaretheorem[
  style=styletheorem,
  name=推论,
  sibling=theorem,
  shaded={bgcolor=\thmshadecolor, margin=0pt},
]{corollary}

\declaretheorem[
  style=styledefinition,
  name=定义,
  sibling=theorem,
  shaded={bgcolor=\defshadecolor, margin=0pt},
]{definition}

\declaretheorem[
  style=styleremark,
  name=注,
  sibling=theorem,
  shaded={bgcolor=\remarkshadecolor, margin=0pt},
]{remark}

\renewcommand{\proofname}{\textcolor{themecolorthm}{证明}}
\renewcommand{\qedsymbol}{\textcolor{themecolorthm}{$\blacksquare$}}

% 彩色标题页（替代 \maketitle）
\makeatletter
\newcommand{\maketitlecolored}{%
  \begin{center}\leavevmode
    {\LARGE\bfseries\color{themeaccent}\@title\par}
    \vspace{0.8em}
    {\large\@author\par}
    {\normalsize\color{black!55}\@date\par}
  \end{center}
  \vspace{1.8em}%
}
\makeatother

% 表格（专业三线表）
\usepackage{booktabs}

% 图表标题样式
\usepackage{caption}
\captionsetup{labelfont={bf,color=themeaccent}, labelsep=period, font=small}

% 链接（置于 preamble 末尾）
\usepackage{hyperref}
\hypersetup{
  colorlinks=true,
  linkcolor=themeaccent!90!black,
  citecolor=themecolordef!75!black,
  urlcolor=themeaccent,
}

% 智能交叉引用（须在 hyperref 之后加载）
\usepackage{cleveref}
\crefname{theorem}{定理}{定理}
\crefname{lemma}{引理}{引理}
\crefname{proposition}{命题}{命题}
\crefname{corollary}{推论}{推论}
\crefname{definition}{定义}{定义}
\crefname{remark}{注}{注}
\crefformat{equation}{式~(#2#1#3)}

\title{GRPO 与 PPO：Rollout、Scoring 与参数更新}
\author{THT}
\date{2026.10.26}

\begin{document}
\maketitlecolored

\section{Rollout：逐个 token 重新生成}

模型根据输入，一个 token 一个 token 地采样，最终拼成一个完整的句子。数学上，第 $t$ 步是在给定输入 $x$ 和已经生成的前缀 $y_{<t}$ 的条件下，从策略分布里抽一个 token：
\[
y_t \sim \pi_\theta(\cdot \mid x,\, y_{<t}),
\qquad
y_{<t} = (y_1, \dots, y_{t-1}).
\]
注意那个 $\sim$：它是\textbf{采样}，不是取最大值。

【举个具体例子】输入 prompt 是「今天天气」。

\begin{enumerate}
  \item 模型跑一次前向，输出词表上的概率分布。假设「不」的概率是 30\%，「很」是 50\%，「真」是 20\%。按概率随机抽样，这次抽中了「很」。
  \item 把「今天天气很」重新输入模型，再跑一次前向，抽出「好」。
  \item 把「今天天气很好」输入模型，抽出「。」，遇到结束符，停止。
\end{enumerate}

于是这一次 Rollout 的结果是「今天天气很好。」。它的关键特征：

\begin{itemize}
  \item \textbf{自回归串行}：必须先生成第 1 个词，才能算第 2 个词的概率。生成 100 个 token 就要跑 100 次前向传播（虽然有 KV Cache 加速，但依然是串行的）。
  \item \textbf{随机性}：涉及采样，同样的输入每次都可能得到不同的句子。
  \item \textbf{极慢}：这正是大模型推理比训练慢得多的原因。在 PPO 里，Rollout 阶段通常占到 80\% 的时间。
\end{itemize}

\section{Scoring：只计算概率}

给定一段完全确定的文本（prompt 加上完整的回答），计算模型生成这段文本的对数概率。数学上就是查表再取对数：
\[
\log \pi_\theta(y_t \mid x,\, y_{<t}).
\]
这里没有采样，只有查表和取对数。

【接着上面的例子】现在手里已经有「今天天气很好。」这句完整的话，把它直接一股脑喂给模型。模型经过一次前向传播（靠因果掩码，也就是三角掩码），可以同时输出所有位置的下一个词的概率。我们不需要关心模型「想」生成什么，只要去看它给出的数字：

\begin{itemize}
  \item 输入「今天天气」后，「很」的概率是 0.6，取对数得 $\log(0.6) \approx -0.51$；
  \item 输入「今天天气很」后，「好」的概率是 0.8，取对数得 $\log(0.8) \approx -0.22$；
  \item 输入「今天天气很好」后，「。」的概率是 0.9，取对数得 $\log(0.9) \approx -0.10$。
\end{itemize}

总的 log 概率是这三项之和：
\[
\log \pi_\theta(y \mid x)
= -0.51 + (-0.22) + (-0.10) = -0.83 .
\]

它的关键特征：

\begin{itemize}
  \item \textbf{完全并行}：整段文本已知，注意力机制可以一次算出所有位置的 log 概率。句子有 100 个词还是 1000 个词，都只跑 1 次前向，而不是像生成那样循环 100 次。
  \item \textbf{确定性}：同样的输入和参数，算出来的概率永远是同一个数。
  \item \textbf{极快}：比 Rollout 快得多。
\end{itemize}

\begin{center}
\begin{tabular}{lll}
\hline
 & Rollout & Scoring \\
\hline
输入 & 截至目前已生成的前缀 & prompt + 完整回答 \\
调用方式 & 逐 token、串行 & 一次前向、并行 \\
输出 & 下一步的概率分布，按概率抽样 & 每个位置的 log 概率 \\
随机性 & 有（采样） & 没有 \\
速度 & 慢，占 PPO 大部分时间 & 快 \\
在 PPO 里做什么 & 产生轨迹（数据） & 算重要性权重 \\
\hline
\end{tabular}
\end{center}

\section{放进 PPO：两个模型、两个阶段}

PPO 的一次迭代里有两个模型同时在场：旧模型 $\pi_{\theta_{\text{old}}}$ 和新模型 $\pi_\theta$。两个阶段各用其中一个。

\subsection{阶段一：Rollout，用旧模型重新生成}

这个阶段用 $\pi_{\theta_{\text{old}}}$ 去生成。生成的目的不是产出最终答案，而是\textbf{拿到数据}：需要完整的句子去让奖励模型打分，去让 critic 估计优势值。生成的同时顺手记下旧模型给出每个词的概率，即
\[
\log \pi_{\theta_{\text{old}}}(y_t \mid x,\, y_{<t}),
\]
这一项之后会当重要性权重的\textbf{分母}。

\subsection{阶段二：Update，用新模型只计算概率}

数据采好之后开始训练新模型 $\pi_\theta$。更新每一次参数时，\textbf{绝对不需要重新生成}，因为数据已经定死了。要做的是：

\begin{enumerate}
  \item 把已有的完整句子喂给 $\pi_\theta$；
  \item 只计算概率，得到新模型给出这些词的
  $\log \pi_\theta(y_t \mid x, y_{<t})$，这是\textbf{分子}；
  \item 算两者的比值（重要性权重）
  \[
  r_t \;=\; \exp\!\left(\log \pi_\theta(y_t \mid x, y_{<t}) - \log \pi_{\theta_{\text{old}}}(y_t \mid x, y_{<t})\right);
  \]
  \item 用 $r_t$ 算 loss，更新参数。
\end{enumerate}

这一步只做前向，不采样，因此可以把同一批数据反复用来更新多次：这正是重要性采样的意义，也是 PPO 比「收集—更新」一次一换的做法更省的地方。

\section{参考}

\begin{enumerate}
\item \url{https://zhuanlan.zhihu.com/p/671780768}
\item \url{https://zhuanlan.zhihu.com/p/1082394115}
\item \url{https://zhuanlan.zhihu.com/p/721073733}
\end{enumerate}

\end{document}
